\documentclass[11pt]{article}
\usepackage{amsmath, amssymb, amsthm}
\usepackage[margin=1in, top=1.2in, bottom=1.2in]{geometry}
\usepackage{parskip}
\usepackage[colorlinks=true, linkcolor=blue, citecolor=blue, urlcolor=blue]{hyperref}
\usepackage{cleveref}
\theoremstyle{definition}
\newtheorem{definition}{Definition}[section]
\newtheorem{theorem}{Theorem}[section]
\newtheorem{lemma}{Lemma}[section]
\newtheorem{proposition}{Proposition}[section]
\newtheorem{corollary}{Corollary}[section]
\theoremstyle{remark}
\newtheorem{remark}{Remark}[section]
\newtheorem{example}{Example}[section]
\setlength{\parindent}{0pt}
\setlength{\parskip}{6pt}
\begin{document}

\begin{center}
{\Large\textbf{1.4 Matrix Group}}\\
\vspace{0.3cm}
{\normalsize\today}
\end{center}
\vspace{0.5cm}


% --- Page 1 ---
Matrix Group

Notation : In $\mathbb{C}^n$, $\overline{e_i} = (0, \cdots, 1^{th}, 0 \cdots)$
is a vector space basis element.

* An endomorphism $\alpha$ of $\mathbb{C}^n$ is defined as
\[ \alpha \overline{e_i} = \sum_{j=1}^n a_{ji} \overline{e_j} \]
where $\alpha = (a_{ij}) \in M_n(\mathbb{C})$.
(We can think of $\alpha \overline{e_i}$ as the $i^{th}$ column
vector.)

We use $\alpha$ for a matrix as well as for the
endomorphism.

We define multiplication of two endomorphisms
$\alpha, \beta$ with corr. matrices $(a_{ij}), (b_{ij})$
and construct $(c_{ij})$ where
\[ c_{ij} = \sum_{k=1}^n a_{ik} b_{kj} \]

% --- Page 2 ---
\begin{lemma}
\label{lem:isomorphism}
$M_n(\mathbb{C}) \cong \mathbb{C}^{n^2}$ as vector spaces.
\end{lemma}

\begin{proof}
Construct the function
\[
f: M_n(\mathbb{C}) \to \mathbb{C}^{n^2}
\]
\[
(a_{ij}) \mapsto (b_{ij})
\]
where $b_{i+(j-1)n} = a_{ij}$.

Observe as $j$ goes from $1$ to $n$, $(j-1)n$ goes from $0$ to $(n-1)n$ and as $i$ goes from $1$ to $n$ as well, we get $i+(j-1)n$ go from $1$ to $n^2$.

Observe $f$ is surjective by construction and if $a_{ij} \neq a_{i'j'}$ then $b_{i+(j-1)n} \neq b_{i'+(j'-1)n}$ hence $f$ is bijective.

Observe $f$ is linear hence $M_n(\mathbb{C}) \cong \mathbb{C}^{n^2}$.
\end{proof}

% --- Page 3 ---
\begin{remark}
Observe we can equip $M_n(\mathbb{C})$ with a topology by ensuring the given map above to be a homeomorphism. By constructing
\[ \mathcal{T}_{M_n(\mathbb{C})} = \{ f^{-1}(CU) \mid U \in \mathcal{T}_{\mathbb{C}^{n^2}} \} \]
\end{remark}

\begin{lemma}
Let $T$ be any topological space,
\[ \phi: T \rightarrow M_n(\mathbb{C}) \]
$\phi$ is continuous if and only if $\pi_{ij} \circ \phi$ is continuous
\begin{proof}
Only if condition is trivial.

If condition:
\[ \pi_{ij} \circ \phi: T \rightarrow \mathbb{C} \] is continuous.

Now $\prod_{1 \leq i,j \leq n} \pi_{ij} \circ \phi: T \rightarrow \mathbb{C}^{n^2}$ is continuous by product topology (or box topology as $\mathbb{C}^{n^2}$ is finite Cartesian product).
\end{proof}
\end{lemma}

% --- Page 4 ---
\[
\text{as } \pi_{ij}^{-1}(U) \text{ for } U \text{ open in } \mathbb{C} \text{ forms a subbasis for product topology on } \mathbb{C}^{n^2}.
\]

\[
\boxed{}
\]

\begin{corollary}
Product of two matrices is continuous
\end{corollary}

\begin{proof}
\[
\text{Pf: } M_n(\mathbb{C}) \times M_n(\mathbb{C}) \xrightarrow{q} M_n(\mathbb{C})
\]

\[
\begin{array}{ccc}
M_n(\mathbb{C}) \times M_n(\mathbb{C}) & \xrightarrow{f \times f} & \mathbb{C}^{n^2} \times \mathbb{C}^{n^2} \\
\downarrow{q} & & \downarrow{p} \\
M_n(\mathbb{C}) & \xrightarrow{f} & \mathbb{C}^{n^2}
\end{array}
\]

where $q: M_n(\mathbb{C}) \times M_n(\mathbb{C}) \rightarrow M_n(\mathbb{C})$,
\[
(\sigma, \tau) \mapsto \sigma \tau
\]
is matrix multiplication. And $p$ is a polynomial
\[
p((b_{ij}), (b'_{ij})) = (c_{ij})
\]
where
\[
c_{ij} = \sum_{k=1}^n b_{ik} b'_{kj}.
\]
Hence $p$ is continuous. Therefore for $q = p \circ (f \times f)$ is continuous. For any open set in $M_n(\mathbb{C})$ $U$, we express
\[
U = f^{-1}(O) \text{ for some open set } O \text{ in }
\]
\end{proof}

% --- Page 5 ---
\[
\mathbb{C}^{n^2}. \text{ Hence } q^{-1}(U) = q^{-1}(f^{-1}(O)) = (f \circ q)^{-1}(O)
\]
which is open as \( f \circ q \) is continuous.

\[
\square
\]

Notation: \( ^t A \) is transpose of a matrix and \( \overline{A} \) is complex conjugate of a matrix.

\begin{lemma}
Transpose and complex conjugate are homeomorphism.
\end{lemma}

\begin{proof}
If \( h: M_n(\mathbb{C}) \rightarrow M_n(\mathbb{C}) \),
\[
\begin{array}{ccc}
A & \xrightarrow{h} & ^t A \\
\end{array}
\]
we can construct \( r: \mathbb{C}^{n^2} \rightarrow \mathbb{C}^{n^2} \) which
is a rearrangement of indices such that
\[
b_{i+j \cdot n} = b_{j+i \cdot n}
\]
which is a homeomorphism.
Hence, \( h = f^{-1} r f \) is a homeomorphism.

If \( k: M_n(\mathbb{C}) \rightarrow M_n(\mathbb{C}) \) is homeomorphic
\[
\begin{array}{ccc}
A & \xrightarrow{k} & \overline{A} \\
\end{array}
\]
\end{proof}

% --- Page 6 ---
as \( k_{ij} \) is homeomorphic for \( 1 \leq i, j \leq n \).

Property of matrices
\[
(\alpha \beta)^t = \beta^t \alpha^t
\]
\[
\frac{d(\alpha \beta)}{dt} = \dot{\alpha} \beta + \alpha \dot{\beta}
\]

* A regular/non singular \( n \times n \) matrix, if it has an inverse i.e. if there exists a matrix \( \sigma^{-1} \) such that \( \sigma \sigma^{-1} = \sigma^{-1} \sigma = E \), where \( E \) is the unit matrix of degree \( n \).

A necessary and sufficient condition for \( \sigma \) to be regular is \( \det \sigma \neq 0 \).

\begin{remark}
If \( \sigma \) is a regular matrix
\begin{enumerate}
\item iff \( \sigma^{-1} \) its reciprocal endomorphism exists
\item then \( ^t (\sigma^{-1}) = (^t \sigma)^{-1} \), \( (\bar{\sigma})^{-1} = (\sigma^{-1})^{-1} \)
\end{enumerate}
\end{remark}

% --- Page 7 ---
\begin{remark}
If $\sigma, \tau$ are regular matrices then $\sigma \tau$ are regular matrices.

$*$ $\det: M_n(\mathbb{C}) \rightarrow \mathbb{C}$ can be expressed as
\[
\begin{array}{ccc}
M_n(\mathbb{C}) & \xrightarrow{\det} & \mathbb{C} \\
\downarrow{f} & & \\
\mathbb{C}^{n^2} & \xrightarrow{p} & \mathbb{C}
\end{array}
\]

$p$ is some polynomial function of $n^2$ variables, which is continuous. We know $f$ is continuous. Hence, $\det$ map is continuous.

\begin{remark}
$GL_n(\mathbb{C}) = \{\sigma \in M_n(\mathbb{C}) \mid \det \sigma \neq 0 \}$
\[
= M_n(\mathbb{C}) \setminus \det^{-1}\{0\}
\]
is a group and an open subset as $\det$ is a continuous map.
\end{remark}
\end{remark}

% --- Page 8 ---
(ii) For $\sigma = (a_{ij}) \in GL_n(\mathbb{C})$ the coefficient
$b_{ij}$ of $\sigma^{-1}$ are given by $b_{ij} = \frac{1}{\det \sigma} A_{ij}$
where $A_{ij}$ is a polynomial of $\sigma$. Hence
mapping $\sigma \mapsto \sigma^{-1}$ is continuous on $GL(n, \mathbb{C})$.
Moreover, it's a homeomorphism.

\begin{remark}
If $\sigma \in GL_n(\mathbb{C})$, we define $\sigma^* = (^t \sigma)^{-1}$.
Then we have $( \sigma \tau )^* = \sigma^* \tau^*$. Hence
$\sigma \mapsto \sigma^*$ is a homeomorphism, and an
automorphism of order 2 of $GL_n(\mathbb{C})$.
\end{remark}

% --- Page 9 ---
\begin{definition}
$\mathrm{O}_n(\mathbb{C}) = \left\{ \sigma \in \mathrm{GL}_n(\mathbb{C}) \mid \sigma^* = \sigma^{-1} \right\}$

is set of orthogonal matrices. If $\sigma = \sigma^*$ then $\sigma$ is complex orthogonal and if $\sigma = \sigma^* = \bar{\sigma}$ then real orthogonal.

$\mathrm{U}_n(\mathbb{C}) = \left\{ \sigma \in \mathrm{GL}_n(\mathbb{C}) \mid \sigma^* = \bar{\sigma}^{-1} \right\}$

is set of all unitary matrices.

Note $\mathrm{U}_n(\mathbb{C})$ or $\mathrm{O}_n(\mathbb{C})$ are non empty as $e \in \mathrm{U}_n(\mathbb{C})$ and $e \in \mathrm{O}_n(\mathbb{C})$.

Notation: $\mathrm{U}_n(\mathbb{C}) = (\mathrm{U}(n))$ and $\mathrm{O}_n(\mathbb{R}) = \mathrm{O}(n)$. Don't be confused!
\end{definition}

% --- Page 10 ---
\begin{lemma}
If $f, g$ are continuous mapping from topological space $X$ to Hausdorff space $Y$ then
\[ \{x \in X \mid f(x) = g(x)\} \] is closed.
\end{lemma}

\begin{proof}
Consider $ \{x \mid f(x) \neq g(x) \} $. Now
\[ x \in \{x \in X \mid f(x) \neq g(x) \} \] implies $f(x) \neq g(x)$
in $Y$. By Hausdorff property there exists open set $U_1$ and $U_2$ such that $f(x) \in U_1$ and $g(x) \in U_2$ and $U_1 \cap U_2 = \emptyset$.
Now $f^{-1}(U_1)$ and $g^{-1}(U_2)$ are open.
Let $V = f^{-1}(U_1) \cap g^{-1}(U_2)$ be a neighborhood of $x$.
For any $y \in V$ $f(y) \neq g(y)$ otherwise
$U_1 \cap U_2 \neq \emptyset$. Hence $x \in V \subseteq \{x \mid f(x) \neq g(x)\}$
is open. Hence $\{x \mid f(x) = g(x)\}$ is closed.
\end{proof}

% --- Page 11 ---
\begin{proposition}
\label{prop:closed_subgroups}
$O_n(\mathbb{C})$, $U(n)$, $O(n)$ are closed subgroups of $GL_n(\mathbb{C})$.
\begin{proof}
Observe mapping $\sigma \mapsto \sigma^*$, $\sigma \mapsto \overline{\sigma}$ and identity mapping are continuous.
By the above lemma, as $\mathbb{R}$ is $T_2$ space.
The sets $U(n) = \{\sigma \mid \overline{\sigma} = \sigma^*\}$ and $O_n(\mathbb{C}) = \{\sigma \mid \sigma = \sigma^*\}$ are closed sets. Observe $GL_n(\mathbb{R}) = \{\sigma \mid \sigma = \overline{\sigma}\}$ is also closed. Hence $U(n)$ and
\[ O(n) = O_n(\mathbb{C}) \cap GL_n(\mathbb{R}) \]
is closed.
As all the given sets are non-empty by subgroup test, we can show these are subgroups.
\end{proof}
\end{proposition}

% --- Page 12 ---
\[
* \quad SL_n(\mathbb{C}) = \left\{ \sigma \in GL_n(\mathbb{C}) \mid \det \sigma = 1 \right\}
\]
is special linear group.

As $\varepsilon \in SL_n(\mathbb{C})$, and by subgroup tests we can show $SL_n(\mathbb{C})$ is a subgroup of $GL_n(\mathbb{C})$. Moreover, $SL_n(\mathbb{C}) = \det^{-1}\{1\}$ implies $SL_n(\mathbb{C})$ is a closed subgroup.

\begin{remark}
\[
SL(n) = GL_n(\mathbb{CIR}) \cap SL_n(\mathbb{C})
\]
\[
SU(n) = SL_n(\mathbb{C}) \cap U(n)
\]
\[
SO(n) = SL_n(\mathbb{C}) \cap O(n)
\]
are closed subgroups of $GL_n(\mathbb{C})$.
\end{remark}

% --- Page 13 ---
\begin{theorem}
\label{thm:compact_groups}
$U(n)$, $O(n)$, $SU(n)$ and $SO(n)$ are compact.
\end{theorem}

\begin{proof}
Observe $O(n)$, $SU(n)$ and $SO(n)$ are closed subgroups of $U(n)$. It is suffice to show $U(n)$ is compact.

Goal: $U(n)$ is homeomorphic to a closed and bounded set of $\mathbb{C}^{n^2}$.

Observe $GL_n(\mathbb{C})$ is a open subgroup of $M_n(\mathbb{C})$ hence it is also a closed subgroup according to a previous lemma. Observe $U(n)$ is closed in $GL_n(\mathbb{C})$ and $GL_n(\mathbb{C})$ is closed in $M_n(\mathbb{C})$ implies $U(n)$ is closed in $GL_n(\mathbb{C})$.

Now for $\sigma = (a_{ij}) \in U(n)$ iff and only if
\[ ^t \overline{\sigma} \sigma = \epsilon \iff \sum_{k} a_{ki} \overline{a_{kj}} = \delta_{ij} \]
\end{proof}

% --- Page 14 ---
implies $\sum_{k=1}^n |a_{ki}|^2 = 1$ which means for $i^{th}$ column vector, the absolute sum is 1.

This implies for any $1 \leq i, j \leq n$ $|a_{ij}| \leq 1$.

Now in $(C^{n^2}, \|\cdot\|_1)$, $v = (v_k)_{k=1}^{n^2}$, $|v_k| \leq 1$ for $1 \leq k \leq n^2$ we define $\|v\|_1 = \max_{1 \leq k \leq n^2} |v_k|$.

Observe $f: M_n(C) \rightarrow C^{n^2}$ as homeomorphism $f(UC_n)$ is closed and bounded, and compact. This implies $UC_n$ is compact.

$\square$

Note: The choice of the norm does not matter as all norms in finite dimension vector space are equivalent (Add the proof to that in the appendix).

\end{document}